% Source: https://artofproblemsolving.com/community/c3h1986142_proportionality_question
% Stable identity: aops-topic-1986142-q01
The sentence means $x=\dfrac{cy}{z}$ for one constant $c$: $x$ is proportional to $y$ when $z$ is held fixed, and inversely proportional to $z$ when $y$ is held fixed. Equivalently $\dfrac{xz}{y}$ is constant. The student's equations would require $\dfrac {x}{y}$ to stay constant even when $z$ changes, and $xz$ to stay constant even when $y$ changes; together they force $y$ to be proportional to $\dfrac{1}{z}$, which is not given. That is the mistake.

From the data, $c=\dfrac{xz}{y}=\dfrac{\dfrac{1}{2}\cdot\dfrac{2}{3}}{\dfrac{3}{4}}=\dfrac{4}{9}$. Then
\[
x=\dfrac{4}{9}\cdot\dfrac{\dfrac{7}{8}}{\dfrac{7}{9}}=\dfrac{4}{9}\cdot\dfrac{9}{8}=\boxed{\dfrac{1}{2}}.
\]
The answer happens to equal the starting value of $x$, which is why the student thought something was wrong.
